% =====================================================================
% Green IT 2026 (MICS2-63): Short Paper, HPC Job Track -- paper v1 (D1)
%
% Built from the professor's template (src/short-paper-hpc.tex of
% student-project-templates.zip) and the Homework 4 copy.
% Every measured figure comes from scripts/v1_numbers.py, which writes
% v1-numbers.tex and v1-tables.tex; rerun it after new runs land in
% results/, then latexmk. Nothing in the text needs retyping.
%
% Build:  latexmk        (run it from the folder above this one)
% =====================================================================
\documentclass[conference,a4paper]{IEEEtran}

\usepackage[T1]{fontenc}
\usepackage{graphicx}
\usepackage{booktabs}
\usepackage{url}
% Let long URLs in the reference list break after slashes, dots and hyphens.
\def\UrlBreaks{\do\/\do\.\do\-\do\_}
\usepackage{amsmath}
% A field of the SLURM job record, or another identifier, set in type.
\newcommand{\field}[1]{\texttt{#1}}
\raggedbottom
% IEEEtran sets a conference abstract in bold; set it in the normal weight.
\usepackage{etoolbox}
\patchcmd{\abstract}{\bfseries}{\mdseries}{}{\PackageWarningNoLine{short-paper}{abstract still bold}}
\usepackage{array}
% A ragged-right paragraph column, for the tables of Appendix A.
\newcolumntype{R}[1]{>{\raggedright\arraybackslash}p{#1}}
% Clickable PDF. The PDF's author field stays empty: the paper is reviewed
% double-blind, and pdfusetitle would copy the author block into it.
\usepackage{xcolor}
\usepackage[colorlinks=true, linkcolor=blue!55!black, citecolor=blue!55!black,
            urlcolor=blue!55!black, bookmarksnumbered=true,
            pdftitle={Life Cycle Assessment of HPCG on the IRIS and Aion Clusters}]{hyperref}
% Generated by scripts/v1_numbers.py; do not edit by hand, rerun the script.
\newcommand{\AionAftEnergy}{27.8}
\newcommand{\AionAftEnergyPct}{$-$30.6\,\%}
\newcommand{\AionAftGrams}{2.51}
\newcommand{\AionAftGramsMax}{2.52}
\newcommand{\AionAftGramsMin}{2.48}
\newcommand{\AionAftGramsSd}{0.009}
\newcommand{\AionAftIters}{200}
\newcommand{\AionAftN}{100}
\newcommand{\AionAftPower}{251}
\newcommand{\AionAftPowerPct}{$-$20.9\,\%}
\newcommand{\AionAftRating}{35.4}
\newcommand{\AionAftRatingPct}{+16.3\,\%}
\newcommand{\AionAftTime}{111.0}
\newcommand{\AionAftTimePct}{$-$12.2\,\%}
\newcommand{\AionAftValid}{all}
\newcommand{\AionAftVsCmp}{$-$20.0\,\%}
\newcommand{\AionAftVsHundred}{$-$20.0\,\%}
\newcommand{\AionAlloc}{54.37\,\%}
\newcommand{\AionBestName}{one process per L3 cache with the 2.0\,GHz cap}
\newcommand{\AionBestSave}{0.63}
\newcommand{\AionBestSavePct}{20.0\,\%}
\newcommand{\AionBusy}{21.8\,\%}
\newcommand{\AionCapA}{2.0}
\newcommand{\AionCapAClock}{2.5}
\newcommand{\AionCapAEmbPct}{+3.4\,\%}
\newcommand{\AionCapAEnergy}{34.1}
\newcommand{\AionCapAEnergyPct}{$-$13.9\,\%}
\newcommand{\AionCapAGrams}{2.97}
\newcommand{\AionCapAGramsPct}{$-$3.9\,\%}
\newcommand{\AionCapAIters}{200}
\newcommand{\AionCapALower}{10 of 10}
\newcommand{\AionCapAN}{10}
\newcommand{\AionCapAOpPct}{$-$13.9\,\%}
\newcommand{\AionCapAPerWorkPct}{$-$3.9\,\%}
\newcommand{\AionCapAPower}{264}
\newcommand{\AionCapAPowerPct}{$-$16.7\,\%}
\newcommand{\AionCapARating}{30.1}
\newcommand{\AionCapARatingPct}{$-$2.5\,\%}
\newcommand{\AionCapARestPct}{+6.0\,\%}
\newcommand{\AionCapASameWork}{yes}
\newcommand{\AionCapASave}{0.12}
\newcommand{\AionCapATime}{129.3}
\newcommand{\AionCapATimePct}{+3.4\,\%}
\newcommand{\AionCapATimed}{76.1}
\newcommand{\AionCapATimedPct}{+1.6\,\%}
\newcommand{\AionCapAValid}{all}
\newcommand{\AionCapB}{1.5}
\newcommand{\AionCapBClock}{2.4}
\newcommand{\AionCapBEmbPct}{+1.8\,\%}
\newcommand{\AionCapBEnergy}{32.4}
\newcommand{\AionCapBEnergyPct}{$-$18.0\,\%}
\newcommand{\AionCapBGrams}{2.89}
\newcommand{\AionCapBGramsPct}{$-$6.6\,\%}
\newcommand{\AionCapBIters}{150/200}
\newcommand{\AionCapBLower}{5 of 10}
\newcommand{\AionCapBN}{10}
\newcommand{\AionCapBOpPct}{$-$18.0\,\%}
\newcommand{\AionCapBPerWorkPct}{+7.8\,\%}
\newcommand{\AionCapBPower}{255}
\newcommand{\AionCapBPowerPct}{$-$19.5\,\%}
\newcommand{\AionCapBRating}{28.4}
\newcommand{\AionCapBRatingPct}{$-$7.7\,\%}
\newcommand{\AionCapBRestPct}{+14.4\,\%}
\newcommand{\AionCapBSameWork}{no}
\newcommand{\AionCapBSave}{0.20}
\newcommand{\AionCapBTime}{127.3}
\newcommand{\AionCapBTimePct}{+1.8\,\%}
\newcommand{\AionCapBTimed}{70.0}
\newcommand{\AionCapBTimedPct}{$-$6.6\,\%}
\newcommand{\AionCapBValid}{all}
\newcommand{\AionCarMetres}{11}
\newcommand{\AionClock}{2.8}
\newcommand{\AionCmpEnergy}{40.1}
\newcommand{\AionCmpGrams}{3.13}
\newcommand{\AionCmpIters}{200}
\newcommand{\AionCmpN}{100}
\newcommand{\AionCmpPower}{317}
\newcommand{\AionCmpRating}{30.5}
\newcommand{\AionCmpTime}{126.5}
\newcommand{\AionCmpWhat}{the 100 runs of the unchanged job in Section~\ref{sec:lcia}}
\newcommand{\AionDiskW}{1.6--5.7}
\newcommand{\AionEmbKg}{1\,200}
\newcommand{\AionEmbOpRatio}{1.4}
\newcommand{\AionEmbOpRatioEU}{0.7}
\newcommand{\AionEmbOpRatioEmber}{1.2}
\newcommand{\AionEnergy}{40.1}
\newcommand{\AionEnergySd}{1.0}
\newcommand{\AionEnergyWh}{11.1}
\newcommand{\AionGEmb}{1.8}
\newcommand{\AionGNode}{1.2}
\newcommand{\AionGOp}{1.3}
\newcommand{\AionGSite}{0.12}
\newcommand{\AionGrams}{3.1}
\newcommand{\AionGramsEU}{4.4}
\newcommand{\AionGramsEmber}{3.3}
\newcommand{\AionGramsNoSite}{3.0}
\newcommand{\AionGramsNodeMax}{3.13}
\newcommand{\AionGramsNodeMin}{3.13}
\newcommand{\AionGramsPueHigh}{3.7}
\newcommand{\AionGramsPueLow}{3.0}
\newcommand{\AionGramsRunMax}{3.25}
\newcommand{\AionGramsRunMin}{3.08}
\newcommand{\AionGramsSd}{0.06}
\newcommand{\AionGramsThree}{3.13}
\newcommand{\AionHot}{69}
\newcommand{\AionHpcgJobs}{3}
\newcommand{\AionHpcgUsers}{1}
\newcommand{\AionHundredGrams}{3.13}
\newcommand{\AionIbIn}{3--7}
\newcommand{\AionIbOut}{20--39}
\newcommand{\AionInstRuns}{100}
\newcommand{\AionLedHours}{5.8}
\newcommand{\AionLife}{4.9}
\newcommand{\AionMgPerGflop}{1.3}
\newcommand{\AionNativeClock}{2.8}
\newcommand{\AionNativeEmbPct}{$-$0.1\,\%}
\newcommand{\AionNativeEnergy}{39.5}
\newcommand{\AionNativeEnergyPct}{$-$0.2\,\%}
\newcommand{\AionNativeGrams}{3.09}
\newcommand{\AionNativeGramsPct}{$-$0.1\,\%}
\newcommand{\AionNativeIters}{200}
\newcommand{\AionNativeLower}{9 of 10}
\newcommand{\AionNativeN}{10}
\newcommand{\AionNativeOpPct}{$-$0.2\,\%}
\newcommand{\AionNativePerWorkPct}{$-$0.1\,\%}
\newcommand{\AionNativePower}{316}
\newcommand{\AionNativePowerPct}{$-$0.1\,\%}
\newcommand{\AionNativeRating}{30.9}
\newcommand{\AionNativeRatingPct}{+0.1\,\%}
\newcommand{\AionNativeRestPct}{$-$0.1\,\%}
\newcommand{\AionNativeSameWork}{yes}
\newcommand{\AionNativeSave}{0.0038}
\newcommand{\AionNativeTime}{124.9}
\newcommand{\AionNativeTimePct}{$-$0.1\,\%}
\newcommand{\AionNativeTimed}{74.9}
\newcommand{\AionNativeTimedPct}{$-$0.1\,\%}
\newcommand{\AionNativeValid}{all}
\newcommand{\AionNodeCount}{4}
\newcommand{\AionNodeGramsMax}{3.2}
\newcommand{\AionNodeGramsMin}{3.1}
\newcommand{\AionNodeSpread}{0\,\%}
\newcommand{\AionNodes}{1}
\newcommand{\AionOptCampaignGrams}{150}
\newcommand{\AionOptRounds}{10}
\newcommand{\AionOptRuns}{50}
\newcommand{\AionOtwoClock}{2.8}
\newcommand{\AionOtwoEmbPct}{$-$0.0\,\%}
\newcommand{\AionOtwoEnergy}{39.6}
\newcommand{\AionOtwoEnergyPct}{+0.2\,\%}
\newcommand{\AionOtwoGrams}{3.10}
\newcommand{\AionOtwoGramsPct}{+0.1\,\%}
\newcommand{\AionOtwoIters}{200}
\newcommand{\AionOtwoLower}{2 of 10}
\newcommand{\AionOtwoN}{10}
\newcommand{\AionOtwoOpPct}{+0.2\,\%}
\newcommand{\AionOtwoPerWorkPct}{+0.1\,\%}
\newcommand{\AionOtwoPower}{317}
\newcommand{\AionOtwoPowerPct}{+0.2\,\%}
\newcommand{\AionOtwoRating}{30.8}
\newcommand{\AionOtwoRatingPct}{+0.0\,\%}
\newcommand{\AionOtwoRestPct}{+0.0\,\%}
\newcommand{\AionOtwoSameWork}{yes}
\newcommand{\AionOtwoSave}{$-$0.0023}
\newcommand{\AionOtwoTime}{125.0}
\newcommand{\AionOtwoTimePct}{$-$0.0\,\%}
\newcommand{\AionOtwoTimed}{74.9}
\newcommand{\AionOtwoTimedPct}{$-$0.1\,\%}
\newcommand{\AionOtwoValid}{all}
\newcommand{\AionOutputKB}{7.15}
\newcommand{\AionPairedLower}{2 of 2}
\newcommand{\AionPairedRounds}{2}
\newcommand{\AionPairedSave}{0.68}
\newcommand{\AionPairedSavePct}{21.7\,\%}
\newcommand{\AionPaybackHi}{2\,438}
\newcommand{\AionPaybackLo}{1\,819}
\newcommand{\AionPaybackYearsHi}{813}
\newcommand{\AionPaybackYearsLo}{606}
\newcommand{\AionPkg}{40.1}
\newcommand{\AionPkgOverTop}{57\,\%}
\newcommand{\AionPkgShareRuns}{100}
\newcommand{\AionPlCapAClock}{2.5}
\newcommand{\AionPlCapAEmbPct}{+3.6\,\%}
\newcommand{\AionPlCapAEnergy}{34.7}
\newcommand{\AionPlCapAEnergyPct}{$-$15.6\,\%}
\newcommand{\AionPlCapAGrams}{3.00}
\newcommand{\AionPlCapAGramsPct}{$-$4.7\,\%}
\newcommand{\AionPlCapAIters}{200}
\newcommand{\AionPlCapALower}{2 of 2}
\newcommand{\AionPlCapAN}{2}
\newcommand{\AionPlCapAOpPct}{$-$15.6\,\%}
\newcommand{\AionPlCapAPerWorkPct}{$-$4.7\,\%}
\newcommand{\AionPlCapAPower}{267}
\newcommand{\AionPlCapAPowerPct}{$-$18.5\,\%}
\newcommand{\AionPlCapARating}{29.9}
\newcommand{\AionPlCapARatingPct}{$-$2.7\,\%}
\newcommand{\AionPlCapARestPct}{+6.2\,\%}
\newcommand{\AionPlCapASameWork}{yes}
\newcommand{\AionPlCapASave}{0.15}
\newcommand{\AionPlCapATime}{129.8}
\newcommand{\AionPlCapATimePct}{+3.6\,\%}
\newcommand{\AionPlCapATimed}{76.4}
\newcommand{\AionPlCapATimedPct}{+1.9\,\%}
\newcommand{\AionPlCapAValid}{all}
\newcommand{\AionPlOptCampaignGrams}{22}
\newcommand{\AionPlOptRounds}{2}
\newcommand{\AionPlOptRuns}{8}
\newcommand{\AionPlRefClock}{2.7}
\newcommand{\AionPlRefEnergy}{41.1}
\newcommand{\AionPlRefGrams}{3.15}
\newcommand{\AionPlRefGramsOpt}{3.15}
\newcommand{\AionPlRefIters}{200}
\newcommand{\AionPlRefN}{2}
\newcommand{\AionPlRefPower}{328}
\newcommand{\AionPlRefRating}{30.8}
\newcommand{\AionPlRefTime}{125.3}
\newcommand{\AionPlRefTimed}{75.0}
\newcommand{\AionPlRefValid}{all}
\newcommand{\AionPlSpreadCapClock}{2.0}
\newcommand{\AionPlSpreadCapEmbPct}{$-$13.0\,\%}
\newcommand{\AionPlSpreadCapEnergy}{27.4}
\newcommand{\AionPlSpreadCapEnergyPct}{$-$33.3\,\%}
\newcommand{\AionPlSpreadCapGrams}{2.46}
\newcommand{\AionPlSpreadCapGramsPct}{$-$21.7\,\%}
\newcommand{\AionPlSpreadCapIters}{200}
\newcommand{\AionPlSpreadCapLower}{2 of 2}
\newcommand{\AionPlSpreadCapN}{2}
\newcommand{\AionPlSpreadCapOpPct}{$-$33.3\,\%}
\newcommand{\AionPlSpreadCapPerWorkPct}{$-$21.7\,\%}
\newcommand{\AionPlSpreadCapPower}{252}
\newcommand{\AionPlSpreadCapPowerPct}{$-$23.4\,\%}
\newcommand{\AionPlSpreadCapRating}{36.1}
\newcommand{\AionPlSpreadCapRatingPct}{+17.4\,\%}
\newcommand{\AionPlSpreadCapRestPct}{$-$8.6\,\%}
\newcommand{\AionPlSpreadCapSameWork}{yes}
\newcommand{\AionPlSpreadCapSave}{0.68}
\newcommand{\AionPlSpreadCapTime}{109.0}
\newcommand{\AionPlSpreadCapTimePct}{$-$13.0\,\%}
\newcommand{\AionPlSpreadCapTimed}{63.1}
\newcommand{\AionPlSpreadCapTimedPct}{$-$15.9\,\%}
\newcommand{\AionPlSpreadCapValid}{all}
\newcommand{\AionPlSpreadClock}{2.7}
\newcommand{\AionPlSpreadEmbPct}{$-$18.2\,\%}
\newcommand{\AionPlSpreadEnergy}{35.0}
\newcommand{\AionPlSpreadEnergyPct}{$-$14.8\,\%}
\newcommand{\AionPlSpreadGrams}{2.62}
\newcommand{\AionPlSpreadGramsPct}{$-$16.7\,\%}
\newcommand{\AionPlSpreadIters}{200}
\newcommand{\AionPlSpreadLower}{2 of 2}
\newcommand{\AionPlSpreadN}{2}
\newcommand{\AionPlSpreadOpPct}{$-$14.8\,\%}
\newcommand{\AionPlSpreadPerWorkPct}{$-$16.7\,\%}
\newcommand{\AionPlSpreadPower}{342}
\newcommand{\AionPlSpreadPowerPct}{+4.2\,\%}
\newcommand{\AionPlSpreadRating}{38.2}
\newcommand{\AionPlSpreadRatingPct}{+24.0\,\%}
\newcommand{\AionPlSpreadRestPct}{$-$16.1\,\%}
\newcommand{\AionPlSpreadSameWork}{yes}
\newcommand{\AionPlSpreadSave}{0.53}
\newcommand{\AionPlSpreadTime}{102.4}
\newcommand{\AionPlSpreadTimePct}{$-$18.2\,\%}
\newcommand{\AionPlSpreadTimed}{60.3}
\newcommand{\AionPlSpreadTimedPct}{$-$19.6\,\%}
\newcommand{\AionPlSpreadValid}{all}
\newcommand{\AionPower}{317}
\newcommand{\AionPue}{1.1}
\newcommand{\AionPueHigh}{1.55}
\newcommand{\AionPueLow}{1.00}
\newcommand{\AionRating}{30.5}
\newcommand{\AionRefClock}{2.8}
\newcommand{\AionRefEnergy}{39.6}
\newcommand{\AionRefGrams}{3.10}
\newcommand{\AionRefGramsOpt}{3.10}
\newcommand{\AionRefIters}{200}
\newcommand{\AionRefN}{10}
\newcommand{\AionRefPower}{316}
\newcommand{\AionRefRating}{30.8}
\newcommand{\AionRefTime}{125.0}
\newcommand{\AionRefTimed}{74.9}
\newcommand{\AionRefValid}{all}
\newcommand{\AionResidentShare}{0.011\,\%}
\newcommand{\AionRuns}{100}
\newcommand{\AionSave}{48\,\%}
\newcommand{\AionSaveMax}{53\,\%}
\newcommand{\AionSaveMin}{34\,\%}
\newcommand{\AionShareEmb}{58\,\%}
\newcommand{\AionShareNode}{39\,\%}
\newcommand{\AionShareOp}{42\,\%}
\newcommand{\AionShareSite}{4\,\%}
\newcommand{\AionTalloc}{2.66}
\newcommand{\AionTdp}{560}
\newcommand{\AionTime}{126}
\newcommand{\AionTimeNodeMax}{126}
\newcommand{\AionTimeNodeMin}{126}
\newcommand{\AionTimeSd}{2}
\newcommand{\AionTopDown}{70.8}
\newcommand{\AionWater}{4.6}
\newcommand{\AionWithinNodeSpread}{6.3\,\%}
\newcommand{\AionWue}{0.413}
\newcommand{\AssistGrams}{430}
\newcommand{\AssistOptGrams}{47}
\newcommand{\AssistOut}{3\,888\,124}
\newcommand{\Gap}{2.8}
\newcommand{\GapEmb}{1.2}
\newcommand{\GapNode}{0.63}
\newcommand{\GapRatio}{1.9}
\newcommand{\GapRatioEU}{2.0}
\newcommand{\GapRatioEmber}{1.9}
\newcommand{\GapRatioPkgOnly}{1.8}
\newcommand{\GapRatioWork}{3.8}
\newcommand{\GapRatioWorkTimed}{3.8}
\newcommand{\GapShareEmb}{41\,\%}
\newcommand{\GapShareNode}{22\,\%}
\newcommand{\GapShareSite}{36\,\%}
\newcommand{\GapShareSiteMax}{46\,\%}
\newcommand{\GapShareSiteMin}{6\,\%}
\newcommand{\GapSite}{1.0}
\newcommand{\Grid}{108.42}
\newcommand{\GridRound}{108}
\newcommand{\IrisAftEnergy}{47.7}
\newcommand{\IrisAftEnergyPct}{$-$18.9\,\%}
\newcommand{\IrisAftGrams}{5.32}
\newcommand{\IrisAftGramsMax}{5.37}
\newcommand{\IrisAftGramsMin}{5.27}
\newcommand{\IrisAftGramsSd}{0.02}
\newcommand{\IrisAftIters}{100}
\newcommand{\IrisAftN}{100}
\newcommand{\IrisAftPower}{220}
\newcommand{\IrisAftPowerPct}{$-$19.9\,\%}
\newcommand{\IrisAftRating}{12.4}
\newcommand{\IrisAftRatingPct}{$-$0.5\,\%}
\newcommand{\IrisAftTime}{216.8}
\newcommand{\IrisAftTimePct}{+1.3\,\%}
\newcommand{\IrisAftValid}{all}
\newcommand{\IrisAftVsCmp}{$-$8.7\,\%}
\newcommand{\IrisAftVsHundred}{$-$11.0\,\%}
\newcommand{\IrisAlloc}{32.94\,\%}
\newcommand{\IrisBestName}{the 2.6\,GHz clock cap}
\newcommand{\IrisBestSave}{0.51}
\newcommand{\IrisBestSavePct}{8.7\,\%}
\newcommand{\IrisBusy}{99.6\,\%}
\newcommand{\IrisCapA}{2.6}
\newcommand{\IrisCapAClock}{2.6}
\newcommand{\IrisCapAEmbPct}{+1.8\,\%}
\newcommand{\IrisCapAEnergy}{47.7}
\newcommand{\IrisCapAEnergyPct}{$-$18.7\,\%}
\newcommand{\IrisCapAGrams}{5.34}
\newcommand{\IrisCapAGramsPct}{$-$8.4\,\%}
\newcommand{\IrisCapAIters}{100}
\newcommand{\IrisCapALower}{9 of 9}
\newcommand{\IrisCapAN}{9}
\newcommand{\IrisCapAOpPct}{$-$18.7\,\%}
\newcommand{\IrisCapAPerWorkPct}{$-$8.4\,\%}
\newcommand{\IrisCapAPower}{219}
\newcommand{\IrisCapAPowerPct}{$-$20.1\,\%}
\newcommand{\IrisCapARating}{12.3}
\newcommand{\IrisCapARatingPct}{$-$0.9\,\%}
\newcommand{\IrisCapARestPct}{+2.6\,\%}
\newcommand{\IrisCapASameWork}{yes}
\newcommand{\IrisCapASave}{0.49}
\newcommand{\IrisCapATime}{217.9}
\newcommand{\IrisCapATimePct}{+1.8\,\%}
\newcommand{\IrisCapATimed}{93.7}
\newcommand{\IrisCapATimedPct}{+0.7\,\%}
\newcommand{\IrisCapAValid}{all}
\newcommand{\IrisCapB}{2.0}
\newcommand{\IrisCarMetres}{21}
\newcommand{\IrisClock}{3.3}
\newcommand{\IrisCmpEnergy}{58.8}
\newcommand{\IrisCmpGrams}{5.82}
\newcommand{\IrisCmpIters}{100}
\newcommand{\IrisCmpN}{9}
\newcommand{\IrisCmpPower}{274}
\newcommand{\IrisCmpRating}{12.4}
\newcommand{\IrisCmpTime}{214.1}
\newcommand{\IrisCmpWhat}{the 9 runs of the unchanged job on the same node}
\newcommand{\IrisDiskW}{2.6--6.7}
\newcommand{\IrisDram}{5.2}
\newcommand{\IrisEmbKg}{1\,256}
\newcommand{\IrisEmbOpRatio}{1.0}
\newcommand{\IrisEmbOpRatioEU}{0.5}
\newcommand{\IrisEmbOpRatioEmber}{0.9}
\newcommand{\IrisEmbSdGrams}{2.9}
\newcommand{\IrisEnergy}{61.0}
\newcommand{\IrisEnergySd}{0.4}
\newcommand{\IrisEnergyWh}{16.9}
\newcommand{\IrisGEmb}{3.0}
\newcommand{\IrisGNode}{1.8}
\newcommand{\IrisGOp}{3.0}
\newcommand{\IrisGSite}{1.2}
\newcommand{\IrisGrams}{6.0}
\newcommand{\IrisGramsEU}{8.8}
\newcommand{\IrisGramsEmber}{6.4}
\newcommand{\IrisGramsNoSite}{4.8}
\newcommand{\IrisGramsNodeMax}{5.98}
\newcommand{\IrisGramsNodeMin}{5.98}
\newcommand{\IrisGramsPkgOnly}{5.7}
\newcommand{\IrisGramsPueHigh}{6.3}
\newcommand{\IrisGramsPueLow}{5.6}
\newcommand{\IrisGramsRunMax}{6.06}
\newcommand{\IrisGramsRunMin}{5.86}
\newcommand{\IrisGramsSd}{0.03}
\newcommand{\IrisGramsThree}{5.98}
\newcommand{\IrisHot}{85}
\newcommand{\IrisHpcgJobs}{7}
\newcommand{\IrisHpcgUsers}{1}
\newcommand{\IrisHundredGrams}{5.98}
\newcommand{\IrisIbIn}{7--10}
\newcommand{\IrisIbOut}{20--40}
\newcommand{\IrisInstRuns}{100}
\newcommand{\IrisLedHours}{11}
\newcommand{\IrisLife}{8.8}
\newcommand{\IrisMemGapShare}{9\,\%}
\newcommand{\IrisMemGrams}{0.25}
\newcommand{\IrisMgPerGflop}{5.1}
\newcommand{\IrisNodeCount}{5}
\newcommand{\IrisNodeGramsMax}{7.2}
\newcommand{\IrisNodeGramsMin}{4.8}
\newcommand{\IrisNodeSpread}{0\,\%}
\newcommand{\IrisNodes}{1}
\newcommand{\IrisOptCampaignGrams}{640}
\newcommand{\IrisOptRounds}{10}
\newcommand{\IrisOptRuns}{120}
\newcommand{\IrisOutputKB}{7.08}
\newcommand{\IrisPairedLower}{9 of 9}
\newcommand{\IrisPairedRounds}{9}
\newcommand{\IrisPairedSave}{0.49}
\newcommand{\IrisPairedSavePct}{8.4\,\%}
\newcommand{\IrisPaybackHi}{3\,010}
\newcommand{\IrisPaybackLo}{2\,246}
\newcommand{\IrisPaybackYearsHi}{430}
\newcommand{\IrisPaybackYearsLo}{321}
\newcommand{\IrisPkg}{55.8}
\newcommand{\IrisPkgOverTop}{92\,\%}
\newcommand{\IrisPkgShareRuns}{100}
\newcommand{\IrisPower}{281}
\newcommand{\IrisPue}{1.63}
\newcommand{\IrisPueHigh}{1.83}
\newcommand{\IrisPueLow}{1.43}
\newcommand{\IrisRating}{12.3}
\newcommand{\IrisRefClock}{3.3}
\newcommand{\IrisRefEnergy}{58.8}
\newcommand{\IrisRefGrams}{5.82}
\newcommand{\IrisRefGramsOpt}{5.82}
\newcommand{\IrisRefIters}{100}
\newcommand{\IrisRefN}{9}
\newcommand{\IrisRefPower}{274}
\newcommand{\IrisRefRating}{12.4}
\newcommand{\IrisRefTime}{214.1}
\newcommand{\IrisRefTimed}{93.1}
\newcommand{\IrisRefValid}{all}
\newcommand{\IrisResidentShare}{0.021\,\%}
\newcommand{\IrisRuns}{100}
\newcommand{\IrisSameNodeVsHundred}{$-$2.5\,\%}
\newcommand{\IrisShareEmb}{50\,\%}
\newcommand{\IrisShareNode}{31\,\%}
\newcommand{\IrisShareOp}{50\,\%}
\newcommand{\IrisShareSite}{19\,\%}
\newcommand{\IrisTalloc}{2.90}
\newcommand{\IrisTdp}{280}
\newcommand{\IrisTime}{217}
\newcommand{\IrisTimeNodeMax}{217}
\newcommand{\IrisTimeNodeMin}{217}
\newcommand{\IrisTimeSd}{1}
\newcommand{\IrisTopDown}{60.8}
\newcommand{\IrisWater}{3.3}
\newcommand{\IrisWithinNodeSpread}{4.5\,\%}
\newcommand{\IrisWue}{0.193}
\newcommand{\LedGrams}{0.54}
\newcommand{\MeasureKg}{0.91}
\newcommand{\OptBestText}{the 2.6\,GHz clock cap cuts a run by 8.7\,\% on IRIS and one process per L3 cache with the 2.0\,GHz cap cuts a run by 20.0\,\% on Aion}
\newcommand{\OptCampaignGramsAll}{1\,100}
\newcommand{\OptCampaignText}{120 runs on IRIS and 158 runs on Aion (plus 13 whose results were not used, in allocations stopped early)}
\newcommand{\OptCapText}{At 2.0\,GHz on Aion, the power changed by $-$16.7\,\% and the run time by +3.4\,\%, so the energy changed by $-$13.9\,\%, the operational emissions by $-$13.9\,\%, the manufacture charged to the run by +3.4\,\%, and the footprint by $-$3.9\,\% (lower in every round). At 1.5\,GHz on Aion, 5 of the 10 runs did 150 iterations instead of 200 and are left out. In the 5 that did the same work, the power changed by $-$19.6\,\% and the run time by +9.7\,\%, so the energy changed by $-$11.8\,\%, the operational emissions by $-$11.8\,\%, the manufacture charged to the run by +9.7\,\%, and the footprint by +0.6\,\% (higher in every round).}
\newcommand{\OptFlagsPctText}{Going back from \texttt{-O3} to \texttt{-O2} changed the footprint of a run by +0.1\,\% on Aion; adding \texttt{-march=native} changed it by $-$0.1\,\% on Aion.}
\newcommand{\OptFlagsText}{On Aion the two levels cannot be told apart from the variation between rounds, and \texttt{-march=native} made no difference beyond that variation.}
\newcommand{\OptLeftOutNote}{ On IRIS, the first of the 10 rounds is left out (Section~\ref{sec:afteropt}).}
\newcommand{\OptPlacementText}{Giving each process its own L3 cache changed the rating by +24.0\,\%, the run time by $-$18.2\,\% and the power by +4.2\,\%, so the footprint changed by $-$16.7\,\% (lower in every round) and the footprint per unit of work by $-$16.7\,\%. With the 2.0\,GHz cap as well, the footprint changed by $-$21.7\,\% (lower in every round). The 2.0\,GHz cap alone changed it by $-$4.7\,\% (lower in every round) in this allocation, against $-$3.9\,\% in the first.}
\newcommand{\OptResultText}{On IRIS, the change kept is the 2.6\,GHz clock cap. Its 100 runs emitted 5.32\,g per run (standard deviation 0.02\,g), against 5.82\,g for the 9 runs of the unchanged job on the same node, a change of $-$8.7\,\%. A run took 216.8\,s instead of 214.1\,s (+1.3\,\%), used 47.7 instead of 58.8\,kJ ($-$18.9\,\%) and rated 12.4\,GFLOP/s, the same to one decimal ($-$0.5\,\%); HPCG validated all of them, each with 100 iterations. The first of the 10 rounds on this node is left out of that comparison: its unchanged run took 245\,s and rated 10.9\,GFLOP/s, against 213 to 215\,s and 12.4 to 12.5\,GFLOP/s for the other 9, and the capped run of that round was slow too (225\,s against 217 to 220\,s); its clock and CPU use were those of the other runs, and nothing else the job recorded explains it. With that round, the change would be $-$9.8\,\%. In the rounds on this node, it emitted less than the unchanged job in 9 of 9 rounds, by 8.4\,\% on average. On another node, the 100 runs of Section~\ref{sec:lcia} emitted 5.98\,g, and this node's unchanged runs 2.5\,\% less: a difference between IRIS's nodes, which is why the comparison rests on the same node. On Aion, the change kept is one process per L3 cache with the 2.0\,GHz cap. Its 100 runs emitted 2.51\,g per run (standard deviation 0.009\,g), against 3.13\,g for the 100 runs of the unchanged job in Section~\ref{sec:lcia}, a change of $-$20.0\,\%. A run took 111.0\,s instead of 126.5\,s ($-$12.2\,\%), used 27.8 instead of 40.1\,kJ ($-$30.6\,\%) and rated 35.4 instead of 30.5\,GFLOP/s (+16.3\,\%); HPCG validated all of them, each with 200 iterations. In the test rounds on one node, it emitted less than the unchanged job in 2 of 2 rounds, by 21.7\,\% on average.}
\newcommand{\OptSaveText}{0.51\,g on IRIS and 0.63\,g on Aion}
\newcommand{\OptValidityText}{On IRIS, HPCG validated all 20 runs of its rounds, all with the same scaled residual. Every configuration did the baseline's 100 iterations. On Aion (flags and clock), HPCG validated all 50 screening runs, all with the same scaled residual. Every configuration did at least the baseline's 200 iterations except the 1.5\,GHz cap, which did 150 in 5 of its 10 runs. HPCG repeats its set of 50 iterations as many times as fit in the 60\,s that \texttt{hpcg.dat} requests, plus one, so a configuration that runs a set slightly slower can do one set fewer, which the rule of Section~\ref{sec:optproc} excludes, and a faster one a set more. On Aion (placement), HPCG validated all 8 screening runs, all with the same scaled residual. Every configuration did the baseline's 200 iterations.}
\newcommand{\PaybackCostHi}{1\,500}
\newcommand{\PaybackCostLo}{1\,100}
\newcommand{\PaybackMeaningText}{The optimisation cost 1\,100 to 1\,500\,g. At the savings measured here, it would take about 1\,819 to 3\,010 comparable runs to recover that cost.}
\newcommand{\PaybackText}{On IRIS the estimated payback is 2\,246 to 3\,010 runs. The 7 matching allocations on IRIS would imply approximately 321 to 430 years under this assumption. On Aion the estimated payback is 1\,819 to 2\,438 runs. The 3 matching allocations on Aion would imply approximately 606 to 813 years under this assumption.}
\newcommand{\PerSecondRatio}{1.1}
\newcommand{\RatingRatio}{2.5}

% Generated by scripts/v1_numbers.py; do not edit by hand, rerun the script.
\newcommand{\RunsTableRows}{%
IRIS & 100 & 217.3 & 61.0 & 55.8 + 5.2 & 281 & 12.3 & 1\,173 & 1.8 & 1.2 & 3.0 & 6.0 \\
Aion & 100 & 126.5 & 40.1 &  & 317 & 30.5 & 2\,345 & 1.2 & 0.12 & 1.8 & 3.1 \\
}
\newcommand{\OtherMeansRows}{%
Timed phase / rest of the run (s) & 94 / 123 & 76 / 51 \\
Peak power over 2\,s (W) & 291 & 322 \\
Node CPU busy & 99.6\,\% & 21.8\,\% \\
Mean core clock (GHz) & 3.3 & 2.8 \\
Processor temperature, mean / hottest (\textdegree C) & 74 / 81 & 61 / 65 \\
Memory touched (GiB) & 27.0 & 27.0 \\
InfiniBand out / in (MB) & 22.5 / 8.5 & 21.5 / 5.0 \\
Ethernet out / in (MB) & 0.16 / 0.19 & 0.15 / 0.13 \\
Local SSD read / written (MB) & 0.00 / 3.0 & 0.01 / 2.1 \\
Idle node power, before and after (W) & 94 & 179 \\
}
\newcommand{\BeforeAfterRows}{%
IRIS & unchanged, Section~\ref{sec:lcia} & 100 & 217.3 & 61.0 & 281 & 12.3 & 100 & 2.99 & 2.98 & 5.98 &  \\
 & unchanged, same node & 9 & 214.1 & 58.8 & 274 & 12.4 & 100 & 2.88 & 2.94 & 5.82 &  \\
 & cap 2.6\,GHz & 100 & 216.8 & 47.7 & 220 & 12.4 & 100 & 2.34 & 2.98 & 5.32 & $-$8.7\,\% \\
\midrule
Aion & unchanged, Section~\ref{sec:lcia} & 100 & 126.5 & 40.1 & 317 & 30.5 & 200 & 1.33 & 1.81 & 3.13 &  \\
 & spread, cap 2.0\,GHz & 100 & 111.0 & 27.8 & 251 & 35.4 & 200 & 0.922 & 1.58 & 2.51 & $-$20.0\,\% \\
}
\newcommand{\OptTableRows}{%
IRIS & unchanged & 9 & 3.3 & 214.1 & 58.8 & 274 & 12.4 & 2.88 & 2.94 & 5.82 &  \\
 & cap 2.6\,GHz & 9 & 2.6 & 217.9 & 47.7 & 219 & 12.3 & 2.34 & 2.99 & 5.34 & $-$8.4\,\% \\
\midrule
Aion (1) & unchanged & 10 & 2.8 & 125.0 & 39.6 & 316 & 30.8 & 1.31 & 1.78 & 3.10 &  \\
 & \texttt{-O2} & 10 & 2.8 & 125.0 & 39.6 & 317 & 30.8 & 1.31 & 1.78 & 3.10 & +0.1\,\% \\
 & \texttt{-march=native} & 10 & 2.8 & 124.9 & 39.5 & 316 & 30.9 & 1.31 & 1.78 & 3.09 & $-$0.1\,\% \\
 & cap 2.0\,GHz & 10 & 2.5 & 129.3 & 34.1 & 264 & 30.1 & 1.13 & 1.84 & 2.97 & $-$3.9\,\% \\
 & cap 1.5\,GHz, 200 it. & 5 & 2.4 & 137.2 & 34.9 & 254 & 28.5 & 1.16 & 1.96 & 3.11 & +0.6\,\% \\
 & cap 1.5\,GHz, 150 it. & 5 & 2.4 & 117.5 & 30.0 & 255 & 28.4 & 0.993 & 1.68 & 2.67 & $-$13.8\,\% \\
\cmidrule(l){2-12}
Aion (2) & unchanged & 2 & 2.7 & 125.3 & 41.1 & 328 & 30.8 & 1.36 & 1.79 & 3.15 &  \\
 & cap 2.0\,GHz & 2 & 2.5 & 129.8 & 34.7 & 267 & 29.9 & 1.15 & 1.85 & 3.00 & $-$4.7\,\% \\
 & spread & 2 & 2.7 & 102.4 & 35.0 & 342 & 38.2 & 1.16 & 1.46 & 2.62 & $-$16.7\,\% \\
 & spread, cap 2.0\,GHz & 2 & 2.0 & 109.0 & 27.4 & 252 & 36.1 & 0.908 & 1.56 & 2.46 & $-$21.7\,\% \\
}

\newcommand{\ProjectRepo}{\url{https://github.com/NexusSin/hpcg-ulhpc-carbon-footprint.git}}

% Guidance notes of the template, kept switched off.
\newif\ifguidance
\guidancefalse
\newcommand{\guidance}[1]{\ifguidance{\itshape\small #1\par\medskip}\fi}

\begin{document}

\title{Life Cycle Assessment of HPCG on the IRIS and Aion Clusters: Environmental Impact and Optimisation}

\author{\IEEEauthorblockN{Fabio Dollaku}
\IEEEauthorblockA{Master in High Performance Computing\\
University of Luxembourg\\
fabio.dollaku.001@student.uni.lu}}

\maketitle

\begin{abstract}
This paper assesses the carbon footprint of the HPCG benchmark on the two
clusters of the University of Luxembourg, air-cooled IRIS and liquid-cooled
Aion, and then tries to lower it without changing the benchmark's code.
Following ISO~14040 and 14044, HPCG was run 100 times on one reserved node of
each cluster with 28 processes. The energy of every run was read from the
processors' counters, and the site's overhead and the node's manufacture were
added. On average a run emits \IrisGrams\,gCO\textsubscript{2}e on IRIS and
\AionGrams\,g on Aion. Embodied and operational emissions stand at
\IrisEmbOpRatio\ to 1 on IRIS and \AionEmbOpRatio\ to 1 on Aion, so the time a
run holds its node matters as much as the power it draws. With the assumed
facility efficiencies, site overhead accounts for \GapShareSite\ of the
estimated gap between the clusters.
HPCG's own build already compiles with \texttt{-O3}. Of the settings a user
controls, compiler flags, a cap on the processor clock and the placement of the
processes, \OptBestText, while
moving the job from IRIS to Aion cuts it by \AionSave.
\end{abstract}

% ---------------------------------------------------------------------

\section{Introduction}\label{sec:intro}

Users of a shared cluster rarely see what their jobs cost the environment. The
scheduler reports time, sometimes energy, but never emissions. The University of
Luxembourg runs two clusters side by side: IRIS, in service since 2017, with
Intel processors cooled by air, and Aion, in service since 2021, with AMD
processors cooled by liquid. A user who sends a job to one of them picks a
processor, a memory system and a cooling technology at once, usually without
noticing.

The job studied here is HPCG, the High Performance Conjugate Gradients
benchmark~\cite{hpcg-repo}. It solves the kind of large sparse linear system found in many
simulation codes, a set of equations in which most coefficients are zero. Its
speed depends more on the memory than on the processor, which makes it a good
test of how a whole node, not just its processor, turns electricity into
computation. The code is HPCG's public reference implementation,
\url{https://github.com/hpcg-benchmark/hpcg}, used unmodified at commit
\texttt{114602d458d1034faa52b71e4c15aba9b3a17698}, with the fixed input of
Table~\ref{tab:frozen}. The method is life cycle assessment, standardised in
ISO~14040~\cite{iso14040} and 14044~\cite{iso14044}. Besides the energy the node draws during the run, it counts the site's
cooling and power distribution, and the run's share of the node's manufacture.

The research question is: how large is the mean carbon footprint of a run of
HPCG on IRIS and on Aion, how much of the difference between the two clusters
comes from the hardware and how much from the cooling, and how much of the
footprint can the settings a cluster user controls remove without touching the
code? The submission scripts, job logs and results are available in the
project repository at \ProjectRepo.

% ---------------------------------------------------------------------
\section{Related Work}

This study draws on four strands of work.

\emph{Estimating the footprint of computation.} Lannelongue et
al.~\cite{lannelongue2021} reduce the footprint of a computation to quantities a
user can supply without measuring anything, such as run time, the cores and
memory used, a usage factor, a power draw per core and per gigabyte, a PUE and a
grid intensity. By their own account the estimate counts the electricity of the
task and leaves out the manufacture of the hardware. Here the energy is measured
on the node and the manufacture is added; on IRIS and Aion it turns out to be
half of the footprint or more (Section~\ref{sec:reporting}).

\emph{Reading energy from the processor.} The inventory rests on the processors'
built-in energy counters. Khan et al.~\cite{khan2018} compared Intel's counters
with power measured at the plug and found the two highly correlated, at a
negligible cost to the running code. They tested Intel processors only; how
accurate the counters are on Aion's AMD processors is still open, and
Section~\ref{sec:threats} counts it among the threats.

\emph{What HPCG measures.} Dongarra, Heroux and Luszczek~\cite{dongarra2016}
introduced HPCG because LINPACK, the benchmark behind the TOP500 list, no longer
reflects the simulations that drive large machines. Those simulations solve
sparse systems, which stress the memory more than the peak arithmetic of the
processor. HPCG reports a rate of floating-point operations and says nothing
about the energy behind it, which is what this paper adds.

\emph{Life-cycle models of computing.} Wadenstein and
Vanderbauwhede~\cite{wadenstein2025} combine the embodied carbon of servers, the
electricity they draw and the site's PUE in one life-cycle model. Where
electricity is low in carbon, they find, the manufacture of the hardware
dominates. Their model covers a computing centre over decades; the same three
terms are applied here to a single job.

% ---------------------------------------------------------------------
\section{Life Cycle Assessment}

The five subsections follow the four phases of ISO~14040~\cite{iso14040} and the
reporting that ISO~14044~\cite{iso14044} requires after them. All tables are in
Appendix~\ref{app:figures}.

\subsection{Goal and scope definition}\label{sec:goalscope}

The goal is to quantify the mean environmental footprint of a run of HPCG,
over 100 runs on each of IRIS and Aion, two clusters that differ in processor
architecture, memory, node design and cooling technology, and to identify the
sources of the difference. Running the same job on both lets us compare the
measured node energy and estimate the contributions of facility overhead and
hardware manufacture under the stated assumptions. The study is aimed at users choosing where
to run jobs and at system operators.

\emph{Functional unit.} One run of HPCG with a fixed input, submitted through
SLURM, the clusters' job scheduler, on one compute node reserved for the job
alone, with the same number of parallel processes (MPI ranks) on both clusters.
Since HPCG adapts its amount of work to a fixed duration, results are reported
per run and per unit of computational work. The job is run 100 times on one node
of each cluster for the baseline. The means and standard deviations describe
the difference between the two nodes and the variation within each set of runs.

\emph{Impact criterion.} The carbon footprint in CO\textsubscript{2}-equivalent
(IPCC~2021~\cite{ipcc2021}, GWP100). Energy, water and hardware use are also
reported in physical units.

\emph{System boundary.} Every part of the cluster that the job uses, in four
components: the compute node, over its whole life cycle; the interconnect and
the shared storage, in their use phase; and the site, comprising power
distribution, cooling and the shared services, such as the login node and the
scheduler, through which the job passes. The user's workstation and the one-off
compilation of the code are excluded.

\subsection{Life cycle inventory}\label{sec:lci}

The inventory quantifies the physical flows of a run, per component and as the
mean of the 100 runs. It rests
on primary data measured on the hardware during the runs, with secondary data
only where a quantity cannot be measured from within a job.

\emph{Measurement protocol.} On each cluster the job runs 100 times back to back
in one allocation, on a node reserved for it alone so that no other workload
enters the readings. No run is discarded. With a hundred runs the mean is steady
enough for the difference between the clusters to stand out from the variation
between runs. The footprint is computed run by run and reported as the mean of
the 100 runs, with its standard deviation. Test jobs in September 2026 established which energy readings each
cluster offers: IRIS's processors report the energy of their packages and of the memory,
Aion's of their packages only. Section~\ref{sec:threats} sizes the difference
this makes.

\emph{Allocation.} Burdens shared by many jobs, such as a node's manufacture, are
allocated to a run in proportion to the time it holds the node, out of all the
time the node is allocated to jobs over its service life.

\subsubsection{Compute node} Processor energy, and memory energy where the
hardware reports it, is read from the processors' built-in energy counters (RAPL)
before and after each run, with the scheduler's energy record as a
cross-check where available. During each run the system also records
the run time, processor utilisation and clock frequency, chip and network-card
temperatures, the memory reserved and used, and local disk traffic; idle power is
measured before and after the runs, and HPCG reports the work done. The
power supply and fans cannot be metered from within a job, a gap stated with the
results. Embodied emissions come from the manufacturer's product carbon
footprint where one exists and from a modelled substitute otherwise, both
secondary sources.

\subsubsection{Interconnect} Traffic through the node's network interfaces is
read from their port counters in every run; the job, confined to one node,
exchanges no data with other compute nodes. The port counters include background
traffic. The measurements provide no per-job estimate of switch electricity.

\subsubsection{Storage} The run directories preserve the benchmark output files.
Their sizes measure the output saved on disk. The scheduler records bytes at
the system-call level, which can include cached reads. These records do not
measure physical traffic or electricity at the shared file servers.

\subsubsection{Site} Cooling and power distribution cannot be metered from
within a job. They enter through the power usage effectiveness (PUE), the ratio
of facility to IT energy. Neither room publishes a measured PUE, so the values
come from European averages for rooms of the same size and from the operator's
design target; water use is estimated from the same averages.

\subsection{Life cycle impact assessment}\label{sec:lcia}

The impact assessment turns the flows of the inventory into an environmental
impact. It uses a single criterion, climate change, characterised with the
IPCC~2021 method~\cite{ipcc2021} and expressed as the global warming potential
over 100 years (GWP100), in grams of CO\textsubscript{2}-equivalent. Every
factor is read in this indicator; a source that counts carbon dioxide alone,
follows an earlier IPCC report or does not say which is flagged where it is
used. A carbon footprint leaves out the other
categories of a full assessment, among them water use, the depletion of
minerals and metals, and toxicity. Water and hardware use still appear in
physical units next to the footprint, but they are not converted into an impact.

\emph{Inventory values.} Table~\ref{tab:runs} summarises the 100 runs on each
cluster. An
average IRIS run holds its node for \IrisTime\,s and draws \IrisEnergy\,kJ
(standard deviation \IrisEnergySd\,kJ) in the processor packages and the memory,
\IrisPower\,W on average. An Aion run takes \AionTime\,s and \AionEnergy\,kJ, or
\AionPower\,W, all of it in the packages, since Aion's processors report no
memory energy. Table~\ref{tab:inventory} gives the mean of every other quantity
measured per run.

\emph{Factors.} Three factors carry the inventory to grams
(Table~\ref{tab:factors}). For the grid I declare \Grid\,gCO\textsubscript{2}
per kWh, the intensity that the Luxembourg regulator, ILR, publishes for all the
electricity supplied to final customers in the country in 2025~\cite{ilr2026},
read on 23 October 2026. Luxembourg imports most of its electricity, so I prefer
a figure for the electricity consumed in the country, imports included, to one
for the electricity produced there. This one counts only the carbon dioxide of
generation, a limit that Section~\ref{sec:threats} discusses, and it applies to
the operational term alone, since the hardware was made in other countries.

For the PUE I declare 1.63 for IRIS, the average that a 2025 assessment for the
European Commission reports for data centres of 100 to 500\,kW~\cite{ec2025},
the class of IRIS's 300\,kW air-cooled room~\cite{ulhpc-dc}. No measured value
is published for Aion either. Its racks are cooled directly by liquid. In every
blade a cold plate, fed by a separate hot-water circuit, takes the place of fans,
and the operator describes the design as targeting a PUE very close to
1~\cite{ulhpc-aion,ulhpc-dc}. I declare 1.1 for Aion, an assumption that
Section~\ref{sec:threats} tests. Applying PUE to the measured processor and
memory energy estimates the associated facility overhead, including cooling.

For manufacture, transport and end of life, an IRIS node carries
\IrisEmbKg\,kg CO\textsubscript{2}e: the 19.3\,\% of Dell's 6\,510\,kg product
carbon footprint for the PowerEdge C6420 that lies outside its use
phase~\cite{dell-c6420-pcf}. The use phase is left out because it is measured. No
such figure exists for the Aion node, so I substitute the \AionEmbKg\,kg that
Wadenstein and Vanderbauwhede model for a recent two-socket AMD
server~\cite{wadenstein2025}. A node's allocated life is its years in service so
far, \IrisLife\ for IRIS's Skylake nodes (since December 2017) and \AionLife\ for
Aion (since late 2021)~\cite{ulhpc-iris,ulhpc-aion}, times the share of each
cluster allocated to jobs over the last twelve months, \IrisAlloc\ and
\AionAlloc\ (Table~\ref{tab:util}). This gives \IrisTalloc\ and \AionTalloc\
years.

\emph{From joules to grams.} The footprint of a run adds an operational and an
embodied term:
\begin{equation*}
C_{\text{run}} = \underbrace{E_{\text{IT}} \cdot \text{PUE} \cdot
I_{\text{grid}}}_{\text{operational}} +
\underbrace{C_{\text{emb}} \cdot \frac{t_{\text{run}}}{T_{\text{alloc}}}}_{\text{embodied}}
\end{equation*}
Here $E_{\text{IT}}$ is the energy measured on the node during the run, converted
from joules to kWh (1\,kWh = 3.6\,MJ), and $I_{\text{grid}}$ is the grid factor
in gCO\textsubscript{2}e per kWh. In the embodied term, $C_{\text{emb}}$ is the
node's embodied footprint in gCO\textsubscript{2}e, $t_{\text{run}}$ the time
the run holds the node, and $T_{\text{alloc}}$ the time the node has been
allocated to jobs over its service life so far, both in seconds. Dividing by the allocated
time instead of the full service life charges the hours a node stands idle to
the jobs that do run. The estimate omits the unmetered IT electricity of
switches and shared storage.

For the mean IRIS run, $E_{\text{IT}}$ is \IrisEnergy\,kJ, or
\IrisEnergyWh\,Wh. At \Grid\,g per kWh this is \IrisGNode\,g for the node's
electricity, and a PUE of 1.63 adds \IrisGSite\,g for cooling and power
distribution. Holding the node for \IrisTime\,s of its \IrisTalloc\ allocated
years carries \IrisGEmb\,g of its manufacture. Averaged over its 100 runs, a run
of HPCG thus emits \IrisGrams\,gCO\textsubscript{2}e on IRIS (standard deviation
\IrisGramsSd\,g), and over the 100 runs on Aion \AionGrams\,g (\AionGNode\,g of
electricity, \AionGSite\,g of site overhead and \AionGEmb\,g of manufacture;
Table~\ref{tab:footprint}). These figures are given to two significant digits,
the precision of the factors.

\subsection{Interpretation}\label{sec:interpretation}

Interpretation, in the sense of ISO~14044, identifies the significant issues,
checks the result for completeness, sensitivity and consistency, and draws
conclusions against the goal of Section~\ref{sec:goalscope}. The figure of
Section~\ref{sec:lcia} is the mean carbon footprint of a run of HPCG with a
fixed input, over 100 runs on one reserved node of each cluster, including the
site's overhead and the run's share of the node's manufacture.

\emph{Embodied versus operational.} The most significant issue is the ratio
between the two terms of the formula. On IRIS a run's share of the node's
manufacture, \IrisGEmb\,g, equals its operational emissions, \IrisGOp\,g of
electricity and site overhead together: a ratio of \IrisEmbOpRatio\ to 1, or
\IrisShareEmb\ embodied against \IrisShareOp\ operational. On Aion the ratio is
\AionEmbOpRatio\ to 1, \AionShareEmb\ against \AionShareOp: its low PUE removes
most of the site overhead, and manufacture becomes the larger part. The ratio is
not a property of the job alone. The grid factor scales only the operational
term, the node's allocated life only the embodied one, and the run time scales
both. At the EU-27 average of 209\,g per kWh~\cite{ries2026s2} the ratio falls
to \IrisEmbOpRatioEU\ to 1 on IRIS and \AionEmbOpRatioEU\ to 1 on Aion, and
operation dominates; on Luxembourg's low-carbon supply, manufacture weighs as
much as electricity or more. A run's length therefore counts as much as its
power, which decides where Section~\ref{sec:optimisation} looks for a saving.

\emph{Meaning.} At the declared grid factor, an hour of a 5\,W LED bulb emits
\LedGrams\,g. An average run on IRIS then amounts to \IrisLedHours\ hours of that bulb,
\IrisCarMetres\,m in a small petrol car at 291\,g per km, or \IrisResidentShare\
of the 28\,kg that a Luxembourg resident emits in a day~\cite{ries2026s2}. An
average run on Aion amounts to \AionLedHours\ hours, \AionCarMetres\,m or
\AionResidentShare.

\emph{Variation between runs.} The 100 runs on each node barely differ. Between
the lowest and the highest, their energy varies by \IrisWithinNodeSpread\ of the
mean on IRIS and by \AionWithinNodeSpread\ on Aion (Table~\ref{tab:runs}). On
Aion, runs 16 to 46 took about 4\,s longer at an unchanged clock and load.
The records do not identify the cause of that slowdown. Those runs stay
in the mean. The first run of each set was no
slower than the rest, so a cold start does not show. The gap between the
clusters is far larger than any of this, so it does not come from chance between
runs.

\emph{Comparing the clusters.} On average a run on IRIS emits \GapRatio\ times
as much as one on Aion, and per unit of work the gap is wider still. Per second of node
time IRIS emits only \PerSecondRatio\ times as much, but Aion's rating is
\RatingRatio\ times higher. Dividing the whole-run footprint by the floating-point
work reported for the timed phase gives \GapRatioWork\ times as much carbon per
GFLOP on IRIS as on Aion (Table~\ref{tab:runs}). Of the estimated gap of
\Gap\,g per run, the model assigns \GapShareSite\ to site overhead at the
declared PUEs, \GapShareNode\ to the node's electricity and \GapShareEmb\ to its
manufacture. Memory is the likeliest
reason why IRIS's run takes longer, which raises both its electricity and its
share of manufacture, although two machines cannot separate memory from
processors (Section~\ref{sec:threats}). HPCG is limited by memory traffic, and an
IRIS node fills only four
of the six memory channels of each processor, for a peak of 154\,GB/s, where an
Aion node fills all sixteen of its channels with faster memory, for
410\,GB/s~\cite{ulhpc-iris,ulhpc-aion} (Table~\ref{tab:frozen}). Aion loses on
another count: the job keeps 28 of its 128 cores busy, \AionBusy\ of the node's
processing time, and the idle cores still draw power.

\emph{Checks.} A top-down estimate tests the counter readings: the processors'
rated power times the time a run holds the node. On IRIS that is
\IrisTdp\,W times the run time, \IrisTopDown\,kJ, and the packages drew
\IrisPkgOverTop\ of it, a plausible level for processors kept fully busy. On
Aion it is \AionTdp\,W times the run time, \AionTopDown\,kJ, of which the
packages drew \AionPkgOverTop, as expected when most cores idle. Neither reading
exceeds its rated value or hints at a fault, and the result keeps the counter
readings. On IRIS the scheduler's own energy record agrees with the counters
within 0.14\,\%. Because the scheduler reads the same registers, the agreement
checks how they are read rather than the sensors themselves. Finally, the sensitivity check of
Section~\ref{sec:threats} recomputes the footprint with other grid factors and
PUE values: the order of the two clusters holds in every case.

\subsection{Reporting}\label{sec:reporting}

Averaged over 100 runs on each cluster, a run of HPCG with the fixed input, on
one node reserved for it with 28 processes, emits \IrisGrams\,gCO\textsubscript{2}e
on IRIS and \AionGrams\,g on Aion (GWP100). The conventions behind
these figures are a grid factor of \Grid\,gCO\textsubscript{2} per kWh, a PUE of
\IrisPue\ for IRIS and \AionPue\ for Aion, and the node's manufacture charged by the time
the run holds the node out of its allocated life. A reader who prefers other
factors can recompute the footprint from the energy per run, \IrisEnergy\,kJ and
\AionEnergy\,kJ at the counters, with the formula of Section~\ref{sec:lcia},
without repeating any measurement.

On both clusters the hotspot is the node's manufacture, with \IrisShareEmb\ of
the IRIS figure and \AionShareEmb\ of Aion's; embodied and operational emissions
stand at \IrisEmbOpRatio\ to 1 on IRIS and \AionEmbOpRatio\ to 1 on Aion. The
node's electricity comes second, with \IrisShareNode\ and \AionShareNode. Both
grow with the time a run holds the node. At the assumed PUEs, site overhead
accounts for \IrisShareSite\ of the total on IRIS and \AionShareSite\ on Aion.

% ---------------------------------------------------------------------
\section{Optimisation}\label{sec:optimisation}

\subsection{Optimisation procedure}\label{sec:optproc}

The hotspot of Section~\ref{sec:reporting} is the node's manufacture, charged by
the time a run holds the node, and the node's electricity comes second. A change
can therefore lower the footprint in two ways: shorten the run, which lowers both
terms, or lower the power drawn without lengthening the run by as much, which
lowers the operational term only. HPCG's code and input stay as they are, and so
do the 28 processes and the node held for the job alone; only what a cluster user
controls changes: how the code is compiled, how fast the processors clock, and
where on the node the processes run.

\emph{Compiler flags.} HPCG's own build configuration, used for every run of
Section~\ref{sec:lcia}, already compiles with \texttt{-O3}, the highest of GCC's
standard optimisation levels, together with \texttt{-ffast-math} and automatic
vectorisation~\cite{gcc-optimize} (Table~\ref{tab:frozen}). Adding
\texttt{-O3} therefore changes nothing. To see what it contributes, the same
build is measured at \texttt{-O2}. The one common flag the configuration leaves
out is \texttt{-march=native}, which lets the compiler use every instruction of
the processor that builds the code~\cite{gcc-x86}. Both were tested on Aion.

\emph{Clock cap.} HPCG spends much of its time waiting for memory
(Section~\ref{sec:interpretation}), so a processor that clocks lower may lose
little speed while drawing less power. SLURM lets a job cap the clock of its
own cores with the \texttt{-{}-cpu-freq} option of \texttt{srun}~\cite{slurm-srun}.
On Aion, whose processors offer only 2.6, 2.0 and 1.5\,GHz steps, the caps tested
are \AionCapA\ and \AionCapB\,GHz. On IRIS, whose cores otherwise run at about
3.3\,GHz with their turbo boost (Table~\ref{tab:inventory}), the cap is
\IrisCapA\,GHz, the processors' base clock, which switches the boost off.

\emph{Placement.} An Aion node has eight memory domains of 16 cores, and each
domain four groups of four cores that share a 16\,MB L3 cache. By default SLURM
puts the job's three or four processes in each domain on the domain's first
group, so they share one cache while three groups stay idle, as the job's record
of each process's cores shows. The spread placement keeps every process in its
domain but gives each its own group, and so its own cache. On IRIS the 28
processes fill the node's 28 cores, and placement has nothing to change.

\emph{Before and after.} The before is the 100 runs of the unchanged job in
Section~\ref{sec:lcia}; the after is 100 runs of the changed job, with the same
protocol, so that the two are measured the same way. On Aion, the change was
chosen in short rounds on one node, each candidate alternating with the
unchanged job in an order that rotated from round to round, so that a slow drift
of the node favours none of them: \AionOptRounds\ rounds for the compiler flags
and clock caps and \AionPlOptRounds\ for the placement, whose difference was
clear from the first round. A change is kept only if HPCG validates every run,
every run does at least the baseline's work, and the footprint is lower in every
round. ``At least'': HPCG fixes its duration, not its work, so a configuration
that runs faster may complete an extra set of iterations, which is extra work at
no extra charge. On IRIS, the \IrisCapA\,GHz cap was chosen without such rounds.
IRIS nodes differ widely, though: in this project's earlier, shorter
measurements on \IrisNodeCount\ IRIS nodes, the same run emitted between
\IrisNodeGramsMin\ and \IrisNodeGramsMax\,g, against \AionNodeGramsMin\ to
\AionNodeGramsMax\,g on \AionNodeCount\ Aion nodes. Section~\ref{sec:lcia}'s IRIS
node was not free in time, so the
job first ran the unchanged and the capped job alternately \IrisOptRounds\ times
on the node it was given, and the 100 capped runs are compared with the
unchanged job on that node. The lecturer suggested the compiler flags; I asked
for the 100 runs after optimisation, for the comparison to rest on them, for a
change that lowers the emissions even if the run gets slower, and for the IRIS
cap without a test. The AI coding agent proposed the \texttt{-O2} control, the
clock caps and the placement, wrote the build variants, the job script and the
script that applies the rule, and submitted and analysed the jobs; I accepted
the levers and the rule for keeping a change.

\subsection{Implementation}\label{sec:optimpl}

Nothing in HPCG's source tree or in \texttt{hpcg.dat} changed. The compiler
variants are passed to \texttt{make} on the command line, as \texttt{CXXFLAGS}
with \texttt{-O2} in place of \texttt{-O3}, or with \texttt{-march=native}
appended, and each is built from the same commit in its own folder at the start
of the allocation. The clock cap is one option of the \texttt{srun} line that
starts HPCG, \texttt{-{}-cpu-freq=<lowest>-<cap>:Performance}, with the
performance governor that the clusters allow; the job records the limit under
which each run actually ran. The placement is another option of the same line,
\texttt{-{}-cpu-bind=map\_cpu:<cores>}, with \texttt{-{}-whole} so that the run may
use any core of the node the job holds, and every run records the cores of its 28
processes. A script applies the rule to the rounds' footprints, computed with the
code that produces every figure of this paper. Two options were rejected. HPCG's build switch
for contiguous matrix storage, \texttt{-DHPCG\_CONTIGUOUS\_ARRAYS}, selects
different code paths in the source, which the brief rules out. Changing the
problem size, the run time or the number of processes would change the
functional unit.

To check that the job still produces the same result, every run's report is
read for HPCG's verdict, the number of conjugate-gradient iterations and the
scaled residual. A configuration that HPCG does not validate, or that does less
work, is no longer the same job.

\subsection{Footprint after optimisation}\label{sec:afteropt}

Table~\ref{tab:beforeafter} compares the 100 runs before and after optimisation;
Table~\ref{tab:opt} gives the rounds that chose the change. Footprints in this
section carry three digits so that the changes show; as absolute figures they
are accurate to two, the precision of the factors (Section~\ref{sec:lcia}).
\OptValidityText

\emph{Compiler flags.} \OptFlagsPctText\ \OptFlagsText

\emph{Clock cap.} \OptCapText

\emph{Placement.} In a second allocation on Aion, with its own baseline:
\OptPlacementText

\emph{The result.} \OptResultText

\subsection{What generalises}\label{sec:generalises}

\emph{The method.} Any job can be examined in four steps, using only the
cluster's own levers. (1)~Split the footprint of one run into its embodied and
operational terms with the formula of Section~\ref{sec:lcia}. Where manufacture
is half of it or more, as on both clusters here, a change that saves energy but
lengthens the run can cost more than it saves. (2)~Choose the node type first:
on the same job, moving from IRIS to Aion saved \AionSave, far more than any
setting tried on one node. (3)~On that node type, try the levers that leave the
code alone, in order of risk: compiler flags passed at build time, then the
placement of the processes and a cap on the clock, both through the scheduler.
(4)~Measure every candidate against the unchanged job in the same allocation,
alternating them, and keep a change only if the program's own check still passes
with at least the same work and if the total, not just the energy, goes down in
every round; then measure the change kept with the full protocol.

\emph{The tool.} Step~(1) is carried out by a script, \texttt{jobprint}, which
needs nothing but a job's identifier. It reads the job's record from the
scheduler, or a saved copy of it when no cluster is at hand, and prints for each
step the time and energy, the footprint split into node, site and manufacture,
the largest of the three, and the reserved cores and memory the job left idle.
Its repository, with a \texttt{README} and a saved record to try it on, is
\ProjectRepo.

% ---------------------------------------------------------------------
\section{Carbon Payback}\label{sec:payback}

The cost of producing the optimisation has two parts. The first is the AI
assistance: over the whole project, \AssistOut\ output tokens, \AssistGrams\,g
CO\textsubscript{2}e at the course factor of 1.03\,Wh per 1\,000 output tokens
and the grid factor of Section~\ref{sec:lcia}; the
session that set up and measured the optimisation accounts for \AssistOptGrams\,g
of it. The second is the cluster time spent testing the configurations and
measuring the kept one: the \OptCampaignText\ emitted \OptCampaignGramsAll\,g
with the factors of Section~\ref{sec:lcia}. The cost is therefore between \PaybackCostLo\,g, counting
that session alone, and \PaybackCostHi\,g, counting all the assistance.

The saving per run is the difference between the 100 runs after optimisation and
the runs they are compared with in Section~\ref{sec:afteropt}, \OptSaveText.
For an illustrative conversion to years, I counted other users' allocations
with \texttt{hpcg} in the job or step name from October 2025 to October 2026.
I assumed one comparable run per matching allocation and the same number of
such allocations each year.
\PaybackText

\subsection{What this means}\label{sec:whatmeans}

In answer to the research question: a run of HPCG emits \IrisGrams\,g
CO\textsubscript{2}e on IRIS and \AionGrams\,g on Aion. The model attributes
\GapShareSite\ of the difference to site overhead. Of the
settings a user controls without touching the code, \OptBestText.

Because manufacture is half of the footprint or more, what counts is how long a
run holds the node. A lower clock acts only on the operational part,
\IrisShareOp\ of a run on IRIS and \AionShareOp\ on Aion, and it lengthens the
run, which charges more of the node's manufacture to it: on Aion the footprint
fell at 2.0\,GHz but rose at 1.5\,GHz, where the runs that did the same work took
longer still. Giving each process its own cache did the opposite: the same work
finished sooner, so both terms fell, and the clock cap on top of it saved
electricity while giving back only part of that time.

\PaybackMeaningText\ Running this job on Aion with the original settings
already reduced its footprint by \AionSave\ relative to IRIS.

% ---------------------------------------------------------------------
\section{Threats to Validity}\label{sec:threats}

The following checks examine how uncertainty in the measurements and model
assumptions affects the comparison.

\emph{What the energy reading covers.} On IRIS the counters cover the processor
packages and the memory, on Aion the packages only. Nothing meters the power
supply, fans, board, network card or local disk, which leaves the operational
term too low on both clusters. Worse, the missing fans skew the comparison.
Air-cooled IRIS nodes need them and liquid-cooled Aion nodes largely do without,
so leaving them out favours IRIS. Reserving the node keeps other jobs out of the reading, but the cores the
job leaves idle are charged to it. That weighs more on Aion, whose nodes have
many more cores than the job uses.

\emph{Two machines, several differences.} IRIS and Aion differ in processor,
memory, node design and cooling at once. The formula estimates a facility
overhead for each system. The comparison cannot isolate the effect of cooling
technology or separate the effect of processor choice from memory configuration. Cooling
acts inside the node as well: a hotter chip leaks more current, and that part of
the cooling's effect is counted as processor energy, outside the PUE. Intel and
AMD implement the counters differently, and the two readings may not be equally
accurate. Aion reports no memory energy, while IRIS's figure includes
\IrisDram\,kJ of it per run: \IrisMemGrams\,g, or \IrisMemGapShare\ of the gap.
On packages alone a run on IRIS would emit \IrisGramsPkgOnly\,g,
\GapRatioPkgOnly\ times Aion's, so the order of the clusters does not depend on
it.

\emph{The declared PUE.} Of all the threats, this one bears most directly on the
research question. Since no measured PUE is published for either room, IRIS's
value comes from European figures for rooms of its size and Aion's from its
operator's design target. With a PUE of \IrisPueLow\ or \IrisPueHigh, a run on
IRIS would emit \IrisGramsPueLow\ or \IrisGramsPueHigh\,g. With \AionPueLow, the
target taken literally, or \AionPueHigh, the European average for rooms of Aion's
size~\cite{ec2025}, a run on Aion would emit \AionGramsPueLow\ or
\AionGramsPueHigh\,g. Across these values the site overhead explains between
\GapShareSiteMin\ and \GapShareSiteMax\ of the gap between the clusters: the
part of the answer that concerns cooling rests on the PUE more than on any
measurement. The
European figure for small rooms comes partly from voluntary reports by more
ambitious operators, which the assessment says may bias it low~\cite{ec2025};
IRIS's true value may be higher. For that reason I also
report the node alone, without any site overhead: \IrisGramsNoSite\,g on IRIS
and \AionGramsNoSite\,g on Aion.

\emph{The grid factor.} Both clusters draw on the same grid. ILR's figure rests on
what the suppliers declare: it counts 52\,\% of the supply as hydropower bought
under contract, and only the carbon dioxide of generation~\cite{ilr2026}. With
Ember's life-cycle intensity of Luxembourg's own production, 123\,g per kWh, a
run would emit \IrisGramsEmber\,g on IRIS and \AionGramsEmber\,g on Aion; with
the EU-27 average of 209\,g~\cite{ries2026s2}, \IrisGramsEU\ and
\AionGramsEU\,g. Ember counts upstream methane with a warming potential of 21, a
value from an earlier IPCC report~\cite{ries2026s2}. A higher factor raises the operational terms only: at the
European average, the node's electricity, not its manufacture, would be the
hotspot on both clusters. The ratio between the clusters barely moves: a run on
IRIS emits \GapRatio\ times as much as one on Aion at the declared factor, and
\GapRatioEU\ times at the European average.

\emph{Substitutes and completeness.} Dell's sheet for the IRIS node states a
standard deviation of $\pm$6\,410\,kg on its 6\,510\,kg, assumes a four-year
life and a configuration with 32\,GB of memory where IRIS nodes hold 128\,GB, and
warns against its use in formal footprints~\cite{dell-c6420-pcf}. Applied to
the manufacture term, that standard deviation alone is $\pm$\IrisEmbSdGrams\,g,
about half of IRIS's total. The substitute
for the Aion node models newer processors and four times its memory, so it
probably overstates its manufacture. Nor are the figures equally complete. The
embodied term covers manufacture, transport and end of life of the compute node.
It excludes manufacture of switches, file servers and cooling equipment. Neither embodied
source states which IPCC report its warming potentials follow. Most of these
gaps make the total too low. The allocation adds its own uncertainty. Charging
manufacture over the years in service so far charges a run the most it can
carry: if IRIS's nodes stay in service for ten years, its manufacture term falls
by an eighth. The allocated share is measured over the whole cluster, where many
failed nodes of IRIS's older generation pull it down; the Skylake nodes alone
are probably busier, which would lower their term further.

\emph{The functional unit.} At a few minutes, the run is shorter and smaller
than the HPCG rules require for an official result, and its figures hold for
this configuration only. Running the same number of processes on both clusters
keeps the job identical but leaves most of an Aion node idle, and a job sized to
fill each node would compare the clusters differently. Each cluster was measured
on one node; another node of the same model may run faster or slower, so the
figures describe the node measured as much as the cluster.

\emph{The optimisation.} On Aion, the 100 runs after optimisation ran on another
node and day than the 100 runs before it, so their difference mixes the change
with the difference between two nodes; Aion's nodes agree within a few per cent,
and the test rounds, which compare the change with the unchanged job on one node,
show how much of it is the change itself. On IRIS, whose nodes differ far more,
the 100 capped runs are compared with \IrisCmpN\ unchanged runs on the same node,
fewer than 100, so that comparison is less precise than the one on Aion. The scheduler caps the
cores' clock, but the memory controller and caches set their own, which a job
can neither see nor set; a cap may act differently on another processor.
\texttt{-march=native} ties the binary to the processor that built it. A single
allocation can contain many runs, as the 100-run baseline jobs illustrate, and
jobs with other names can also launch HPCG. The accounting records therefore
leave the actual annual number of comparable runs unknown.

\emph{Burden shifting and the single criterion.} A carbon-only result hides
whatever it does not count. Liquid cooling may change how much water a site
uses: at the European averages for the two size classes, 0.193 and 0.413\,L per
kWh of IT energy~\cite{ec2025}, a run consumes about \IrisWater\,mL of water on
IRIS and \AionWater\,mL on Aion, although direct liquid cooling can need little
water, so Aion's figure may be too high. Replacing an air-cooled cluster with a
liquid-cooled one means new hardware to make and old hardware to dispose of. The
results cover climate
change only, for one job on two clusters at one site, and support no comparison
with other jobs.

% ---------------------------------------------------------------------
% Unnumbered on purpose: a disclosure statement sits outside the argument.
\section*{AI Usage Transparency}

I used AI to analyse the architecture of the IRIS and Aion clusters, adapt the
SLURM submission script, and interpret the logs from the HPCG runs. It also
helped me extract information from the cited websites and apply findings from
the literature to this job.

% ---------------------------------------------------------------------
\bibliographystyle{IEEEtran}
\bibliography{refs}

% ---------------------------------------------------------------------

\appendices
\section{Figures and Tables}\label{app:figures}

\begin{table*}[!t]
\caption{The job and the two machines it ran on}
\label{tab:frozen}
\centering
\footnotesize
\begin{tabular}{@{}R{0.12\textwidth}R{0.41\textwidth}R{0.41\textwidth}@{}}
\toprule
 & IRIS & Aion \\
\midrule
Code & \multicolumn{2}{>{\raggedright\arraybackslash}p{0.83\textwidth}}{HPCG 3.1,
\url{https://github.com/hpcg-benchmark/hpcg}~\cite{hpcg-repo}, commit
\texttt{114602d458d1034faa52b71e4c15aba9b3a17698} of 2023-09-23, reference
kernels, unmodified} \\
Build & \multicolumn{2}{>{\raggedright\arraybackslash}p{0.83\textwidth}}{HPCG's
\texttt{MPI\_GCC\_OMP} setup: \texttt{mpicxx -O3 -ffast-math -ftree-vectorize
-fopenmp}; GCC 13.2.0 and Open MPI 4.1.6, built inside each allocation} \\
Input & \multicolumn{2}{>{\raggedright\arraybackslash}p{0.83\textwidth}}{\texttt{hpcg.dat}, 84 bytes, SHA-256
\texttt{8a18c3bde4485866bb24d3290da206003a4690112f2b01d6e77d9a3f7913376b}; line 3
\texttt{104 104 104} (grid points per process), line 4 \texttt{60} (target
seconds of the timed phase)} \\
Job & \multicolumn{2}{>{\raggedright\arraybackslash}p{0.83\textwidth}}{28 MPI processes of one thread on one
exclusive node, bound to cores; 100 runs back to back in one allocation per
cluster, October 2026} \\
Partition & \texttt{batch}, Skylake nodes & \texttt{batch} \\
Node & Dell PowerEdge C6420 & one of the three nodes of an Atos BullSequana X2410
blade \\
Processors & 2 $\times$ Intel Xeon Gold 6132: 14 cores, 2.6\,GHz, 140\,W & 2
$\times$ AMD EPYC 7H12: 64 cores, 2.6\,GHz, 280\,W \\
Memory & 128\,GB DDR4-2400, 4 of 6 channels per processor, peak 154\,GB/s &
256\,GB DDR4-3200, all 8 channels per processor, peak 410\,GB/s \\
Cooling, room & air, in-row units; 300\,kW room & direct liquid, hot-water loop;
1\,050\,kW room \\
In service & Skylake nodes since December 2017 & since late 2021 \\
Energy read & RAPL counters: packages and memory & RAPL counters: packages only \\
\bottomrule
\end{tabular}
\par\smallskip
\parbox{\textwidth}{\footnotesize Hardware, rooms and dates from the ULHPC
documentation~\cite{ulhpc-iris,ulhpc-aion,ulhpc-dc}; versions and checksum from
the job logs. Peak memory bandwidth = channels $\times$ transfer rate $\times$
8 bytes.}
\end{table*}

\begin{table}[!t]
\caption{Declared factors}
\label{tab:factors}
\centering
\footnotesize
\setlength{\tabcolsep}{3pt}
\begin{tabular}{@{}R{0.36\columnwidth}R{0.25\columnwidth}R{0.33\columnwidth}@{}}
\toprule
Factor & Value & Source, basis \\
\midrule
Grid intensity & \Grid\,gCO\textsubscript{2}/kWh & ILR, supply to final customers, 2025, CO\textsubscript{2} only~\cite{ilr2026} \\
PUE, IRIS & \IrisPue & EU, 100--500\,kW class~\cite{ec2025} \\
PUE, Aion & \AionPue & assumption: the operator's target for its liquid cooling, ``very close to 1''~\cite{ulhpc-aion} \\
Node manufacture, IRIS & \IrisEmbKg\,kg CO\textsubscript{2}e & Dell sheet, share outside use~\cite{dell-c6420-pcf} \\
Node manufacture, Aion & \AionEmbKg\,kg CO\textsubscript{2}e & modelled AMD node, substitute~\cite{wadenstein2025} \\
Allocated life, IRIS & \IrisLife\,y $\times$ \IrisAlloc\ = \IrisTalloc\,y & ULHPC timeline; \texttt{sreport} \\
Allocated life, Aion & \AionLife\,y $\times$ \AionAlloc\ = \AionTalloc\,y & ULHPC timeline; \texttt{sreport} \\
AI assistance & 1.03\,Wh per 1\,000 output tokens & course factor \\
\bottomrule
\end{tabular}
\par\smallskip
\parbox{\columnwidth}{\footnotesize Read 29 September and 23 October 2026.
Comparators: 5\,W LED at the grid factor, 291\,g per km in a small petrol car,
28\,kg per Luxembourg resident per day~\cite{ries2026s2}.}
\end{table}

\begin{table*}[!t]
\caption{Mean per run over the 100 runs on each cluster}
\label{tab:runs}
\centering
\footnotesize
\setlength{\tabcolsep}{4pt}
\begin{tabular}{@{}lrrrlrrrrrrr@{}}
\toprule
Cluster & Runs & Time & Energy & Packages + memory & Power & Rating & Timed work &
\multicolumn{4}{c}{gCO\textsubscript{2}e per run} \\
\cmidrule(l){9-12}
 & & (s) & (kJ) & (kJ) & (W) & (GFLOP/s) & (GFLOP) & node & site & manufacture & total \\
\midrule
\RunsTableRows
\bottomrule
\end{tabular}
\par\smallskip
\parbox{\textwidth}{\footnotesize Every value is the mean of the 100 runs on one
node of each cluster. The footprint of a single run has a standard deviation of
\IrisGramsSd\,g on IRIS and \AionGramsSd\,g on Aion, and ranges from
\IrisGramsRunMin\ to \IrisGramsRunMax\,g and from \AionGramsRunMin\ to
\AionGramsRunMax\,g. Energy: RAPL counters read between runs. Timed work: the
floating-point operations of HPCG's timed phase, 100 CG iterations per run on
IRIS and 200 on Aion. Grams with the factors of Table~\ref{tab:factors}.}
\end{table*}

\begin{table}[!t]
\caption{Mean footprint of a run, over 100 runs per cluster}
\label{tab:footprint}
\centering
\footnotesize
\setlength{\tabcolsep}{4pt}
\begin{tabular}{@{}lrr@{}}
\toprule
gCO\textsubscript{2}e per run & IRIS & Aion \\
\midrule
Node electricity & \IrisGNode\ (\IrisShareNode) & \AionGNode\ (\AionShareNode) \\
Site overhead (PUE) & \IrisGSite\ (\IrisShareSite) & \AionGSite\ (\AionShareSite) \\
Node manufacture & \IrisGEmb\ (\IrisShareEmb) & \AionGEmb\ (\AionShareEmb) \\
\midrule
Total & \IrisGrams & \AionGrams \\
Standard deviation & \IrisGramsSd & \AionGramsSd \\
Embodied : operational & \IrisEmbOpRatio\ : 1 & \AionEmbOpRatio\ : 1 \\
\midrule
Hours of a 5\,W LED & \IrisLedHours & \AionLedHours \\
Metres in a small petrol car & \IrisCarMetres & \AionCarMetres \\
Share of a resident's day & \IrisResidentShare & \AionResidentShare \\
\midrule
At 123\,g/kWh (Ember, LU) & \IrisGramsEmber & \AionGramsEmber \\
At 209\,g/kWh (EU-27) & \IrisGramsEU & \AionGramsEU \\
\quad embodied : operational & \IrisEmbOpRatioEU\ : 1 & \AionEmbOpRatioEU\ : 1 \\
At PUE \IrisPueLow\ / \IrisPueHigh\ (IRIS) & \IrisGramsPueLow\ / \IrisGramsPueHigh & \\
At PUE \AionPueLow\ / \AionPueHigh\ (Aion) & & \AionGramsPueLow\ / \AionGramsPueHigh \\
Node alone, no site overhead & \IrisGramsNoSite & \AionGramsNoSite \\
\bottomrule
\end{tabular}
\par\smallskip
\parbox{\columnwidth}{\footnotesize \IrisRuns\ runs on IRIS, \AionRuns\ on Aion.
Shares of the total in brackets. Operational = node electricity + site overhead.}
\end{table}

\begin{table}[!t]
\caption{Other measured quantities: mean per run over the 100 runs}
\label{tab:inventory}
\centering
\footnotesize
\setlength{\tabcolsep}{3pt}
\begin{tabular}{@{}R{0.46\columnwidth}R{0.24\columnwidth}R{0.24\columnwidth}@{}}
\toprule
Quantity per run & IRIS & Aion \\
\midrule
Cores held / used & 28 / 28 & 128 / 28 \\
\OtherMeansRows
Memory reserved (GiB) & 112 & 219 \\
Data between compute nodes & none, one node & none, one node \\
Saved HPCG output (kB/run) & \IrisOutputKB & \AionOutputKB \\
Water, declared & \IrisWater\,mL & \AionWater\,mL \\
\bottomrule
\end{tabular}
\par\smallskip
\parbox{\columnwidth}{\footnotesize Means of the 100 runs on each cluster, the
node sampled every 2\,s. Power: packages and memory on IRIS, packages only on
Aion; idle power over 60\,s before and after the runs. Water: litres per kWh of IT
energy for each room's size class, 0.193 and 0.413~\cite{ec2025}.}
\end{table}

\begin{table}[!t]
\caption{Utilisation of the two clusters over twelve months}
\label{tab:util}
\centering
\footnotesize
\begin{tabular}{@{}lrr@{}}
\toprule
Share of capacity & IRIS & Aion \\
\midrule
Allocated, to jobs or reservations & 32.94\,\% & 54.37\,\% \\
Idle & 29.06\,\% & 35.86\,\% \\
Down & 19.09\,\% & 2.67\,\% \\
Planned & 16.50\,\% & 5.64\,\% \\
Planned down & 2.41\,\% & 1.47\,\% \\
\bottomrule
\end{tabular}
\par\smallskip
\parbox{\columnwidth}{\footnotesize \texttt{sreport cluster utilization},
2025-09-29 to 2026-09-28, whole clusters. Many of IRIS's older nodes had failed,
so its share probably understates the Skylake nodes' own.}
\end{table}

\begin{table*}[!t]
\caption{Before and after optimisation: the unchanged job and the change kept, mean per run}
\label{tab:beforeafter}
\centering
\footnotesize
\setlength{\tabcolsep}{4pt}
\begin{tabular}{@{}llrrrrrrrrrr@{}}
\toprule
Cluster & Version & Runs & Time & Energy & Power & Rating & Iterations &
\multicolumn{4}{c}{gCO\textsubscript{2}e per run} \\
\cmidrule(l){9-12}
 & & & (s) & (kJ) & (W) & (GFLOP/s) & & operational & manufacture & total & change \\
\midrule
\BeforeAfterRows
\bottomrule
\end{tabular}
\par\smallskip
\parbox{\textwidth}{\footnotesize Same code, input, 28 processes, node held for
the job alone, protocol and factors as Table~\ref{tab:runs}. Change: against the
100 runs of Section~\ref{sec:lcia} on Aion, and against the unchanged runs on the
same node on IRIS, whose nodes differ widely (Section~\ref{sec:optproc}).\OptLeftOutNote\ Spread:
one process per L3 cache. HPCG validated every run. gCO\textsubscript{2}e to
three digits so that the change shows; accurate to two (Section~\ref{sec:lcia}).}
\end{table*}

\begin{table*}[!t]
\caption{The rounds that chose the change: every configuration against the unchanged job on the same node}
\label{tab:opt}
\centering
\footnotesize
\setlength{\tabcolsep}{4pt}
\begin{tabular}{@{}llrrrrrrrrrr@{}}
\toprule
Cluster & Configuration & Runs & Clock & Time & Energy & Power & Rating &
\multicolumn{4}{c}{gCO\textsubscript{2}e per run} \\
\cmidrule(l){9-12}
 & & & (GHz) & (s) & (kJ) & (W) & (GFLOP/s) & operational & manufacture & total & change \\
\midrule
\OptTableRows
\bottomrule
\end{tabular}
\par\smallskip
\parbox{\textwidth}{\footnotesize Means over the rounds of one allocation, on one
node, each round running every configuration once in a rotating order. Unchanged:
HPCG's own build (\texttt{-O3}) at the default clock and placement, as in
Section~\ref{sec:lci}. Clock: mean over all the node's cores, sampled every
2\,s; on Aion it averages the 28 cores that run HPCG with the 100 idle ones, so
it stays above the cap, which the job recorded on the cores. Same code, input, 28
processes and factors as Table~\ref{tab:runs}; change relative to the unchanged
job of the same allocation. (1), (2): Aion's two allocations; the second stopped
after two rounds, once the result was clear. Spread: one process per L3 cache
(Section~\ref{sec:optproc}); with the cap, it applied to all the node's cores.
On IRIS, the cap was chosen beforehand, and its rounds give the unchanged job on
the node of the 100 capped runs.\OptLeftOutNote\ HPCG validated every run; ``it.'': a
configuration whose runs did different numbers of iterations is split by that
number (Section~\ref{sec:afteropt}). gCO\textsubscript{2}e to three digits,
accurate to two.}
\end{table*}

\clearpage
\section{AI Usage Journal}\label{app:journal}

I used AI to help me read the technical documentation for IRIS and Aion and
understand how the differences between their nodes affect this particular HPCG
job. One question was what it meant to request 28 cores on each machine. On an
IRIS node, the 28 processes occupy all the cores across its two Intel processors.
On an Aion node, those same processes occupy 28 of the 128 cores across its two
AMD processors. AI helped me interpret the specifications for memory channels
and shared caches in the cluster documentation~\cite{ulhpc-iris,ulhpc-aion}.
These details matter for HPCG because its
performance depends on how quickly each process can access the data it needs.
I also used AI to adapt the SLURM submission script to the two clusters. This
included understanding how the script requests an entire node, assigns one
thread to each MPI process and binds processes to cores. The script gives
each run its own directory and copies the same input file into it, so I could
trace a benchmark report back to the run that produced it. Keeping the input and
process count the same gave me a consistent basis for comparing the runs.
Reserving the node for the job also meant that another user's workload would
not contribute to the energy readings. Once the jobs had finished, I used AI
to help me read the scheduler logs, the energy records and the reports produced
by HPCG. Each source answered a different question about a run. The job record
identified the node and allocation, the counter readings taken before and after
the run gave the energy consumed by the components being monitored, and the
benchmark report recorded the work completed and the performance achieved.
Reading them together helped me relate a difference in elapsed time to the
work done. It also helped me interpret a gap in the energy data. IRIS reports
energy for both the processor packages and memory. Aion reports package energy
only, so memory consumption on Aion is an unknown in the comparison. I also
used AI to help extract information from the cited websites and reports and
work out which published values applied to this job. For the electricity
calculation, the Luxembourg regulator reports a factor for electricity supplied
to customers, including imports~\cite{ilr2026}. That figure applies to the
electricity used while the job runs. Dell's product footprint covers a server's
whole life~\cite{dell-c6420-pcf}. Using its full value would count the use phase twice, so the
calculation uses the share outside that phase and allocates a fraction to the
time the job occupied the node.

\end{document}
